\section{Finite Groups}


\subsection{Order and Subgroups}

\begin{defbox}
   Let $G$ be a finite group. We call the \textbf{order} of $G$, denoted $|G|$, to be its cardinality.
\end{defbox}

\begin{example}
   The order of the additive group $\ZZ / n\ZZ$ is $n$ while the order of the multiplicative group $(\ZZ / n\ZZ)^*$ is eqaul to Euler's totient function $\phi(n)$.
\end{example}

\begin{example}
   For groups $(G, \cdot_G)$ and $(H, \cdot_H)$, we can form their \textit{direct product} as the group with the underlying set $G \times H$ and binary operation $\cdot \colon (G \times H)^2 \to G \times H$ given by coordinatewise multiplication: 
   \[ (g_0, h_0) \cdot (g_1, h_1) := (g_0 \cdot g_1, h_0 \cdot h_1) \]
We can easily check that this is a group with identity $(e_G, e_H)$ and inverse $(g^{-1}, h^{-1})$. 

For example, the product of $(\ZZ, +)$ with itself is just $(\ZZ^2, +)$, and the same idea is true for the product of $(\RR, +)$ with itself. If $G$ and $H$ are finite groups, then the order of $G \times H$ is $|G||H|$. 
\end{example}

\begin{defbox}
   Let $(G, \cdot)$ be a group. We say that a \textbf{subgroup} of $G$ is a subset $H \subseteq G$ which is a group of its own with the same operations, i.e. $H$ satisfies:
   \begin{enumerate}
      \item $e_G \in H$. 
      \item For all $g_0$, $g_1 \in G$, if $g_0$, $g_1 \in H$, then $g_0 \cdot g_1 \in H$.
      \item For all $g \in G$, if $g \in H$ then $g^{-1} \in H$. 
   \end{enumerate}
   If $H$ is a subgroup of $G$, we write $H \leq G$. 
\end{defbox}

\begin{remark}
   For any group $G$, the singleton $\{e_G\}$ is its trivial subgroup. Also $G \leq G$.
\end{remark}

\begin{example}
   For $(\ZZ, +)$, $2\ZZ := \{2n \colon n \in \ZZ\}$ is a subgroup of $(\ZZ, +)$. Indeed, $0 \in 2\ZZ$, for $m$, $n \in 2\ZZ$, $m + n \in 2\ZZ$, and if $n \in 2\ZZ$, then $-n \in 2\ZZ$. Indeed, any $d\ZZ := \{dn \colon d \in \ZZ, n \in \ZZ\}$ is a subgroup of $(Z, +)$. However, $2\ZZ + 1 \not\leq \ZZ$, as $0 \notin 2\ZZ + 1$. 
\end{example}

\begin{example}
   For the multiplicative group $(\CC^*, \cdot)$, where $\CC^* := \CC \setminus \{0\}$, the following are all subgroups:
   \begin{enumerate}
      \item $\QQ^*$ and $\RR^*$;
      \item $\RR^* \cup i\RR^*$;
      \item $\{\pm1\}$, $\{\pm 1, \pm i\}$; 
      \item The unit circle $S := \{z \in \CC : |z| = 1\}$. Indeed, $1 \in S$, if $z_0$, $z_1 \in S$, then $|z_0z_1| = |z_0|\cdot|z_1| = 1 \cdot 1 = 1$, and $|z_0^{-1}| = |z_0|^{-1} = 1^{-1} = 1$. 
   \end{enumerate}
\end{example}

\subsection{Cyclic Groups}

\begin{defbox}
   Let $G$ be a group and $g \in G$. We call the set $\langle g \rangle := \{g^n : n \in \ZZ\}$ the \textbf{subgroup of $G$ generated by $g$}. Indeed $\langle g \rangle$ is easily verified to be a subgroup of $G$ because $g^0 = e_G$, $g^n \cdot g^m = g^{n + m}$ and $(g^n)^{-1} = g^{-n}$. The order of this group is called the \textbf{order} of $g$ and denoted $|\langle g \rangle|$.  
\end{defbox}

\begin{lemmabox}
   Let $G$ be a group and $g \in G$. Then $|\langle g \rangle|$ is equal to the least $k \in \NN^+$ such that $g^k = e_G$. 
\end{lemmabox}